\chapter{三角 Fourier 级数}
\label{chap:trigonometric-fourier-series}

周期函数可以用三角多项式逼近，并不意味着它的 Fourier 部分和必然收敛。
这两个问题分别由 Dirichlet 核和 Fejér 核控制：前者有剧烈振荡且
\(L^1\) 范数随阶数增长，后者非负、总质量固定并集中到原点。本章先在点态
意义下考察这一区别。所有函数都视为 \(2\pi\)-周期函数，积分取任意一个
长度为 \(2\pi\) 的周期；记圆周为 \(\mathbb T=\R/(2\pi\mathbb Z)\)。

本章的 \(L^p\) 范数使用一个周期上的长度测度，卷积公式中的
\(1/(2\pi)\) 则显式保留。以后使用归一化测度时会另作说明。
连续周期函数在 \(L^p(\mathbb T)\) 中稠密（\(1\le p<\infty\)）：
先将 \(f\) 在区间端点的小邻域截去，所产生的 \(L^p\) 误差趋零；
再用\theoremref{thm:cc-dense-lp}在实线上逼近这个零延拓函数。
乘以一个支撑于 \((-\pi,\pi)\)、在截断函数支撑上为一且取值于
\([0,1]\) 的连续截断函数，不会增加误差。所得函数可连续地周期延拓。
平移保持 \(L^p\) 范数，因此连续周期函数的一致连续性和上述稠密性
也给出圆周上的 \(L^p\) 平移连续性。

\section{系数与部分和}

在 \(L^1(\mathbb T)\) 中定义
\[
 \widehat f(k)=\frac1{2\pi}\int_{-\pi}^{\pi}f(x)e^{-ikx}\,dx,
 \qquad k\in\mathbb Z.
\]
若 \(f\) 实值，则 \(\widehat f(-k)=\overline{\widehat f(k)}\)，从而复指数
形式与通常的正弦、余弦形式相互换算。第 \(N\) 个对称部分和为
\[
 S_Nf(x)=\sum_{k=-N}^{N}\widehat f(k)e^{ikx}.
\]
有限求和允许与积分换序，所以
\begin{equation}\label{eq:fourier-dirichlet-convolution}
 S_Nf(x)=\frac1{2\pi}\int_{-\pi}^{\pi}f(x-t)D_N(t)\,dt,
 \qquad D_N(t)=\sum_{k=-N}^Ne^{ikt}.
\end{equation}
对 \(t\notin2\pi\mathbb Z\) 求有限等比和，得
\[
 D_N(t)=\frac{\sin((N+\tfrac12)t)}{\sin(t/2)},\qquad
 \frac1{2\pi}\int_{-\pi}^{\pi}D_N(t)\,dt=1.
\]
因此 \(D_N\) 具有正确总质量，却不是非负核。下面核实其绝对积分的增长
及这一增长对部分和算子的影响。

令 \(\lambda=N+\tfrac12\)。由
\(t/\pi\le\sin(t/2)\le t/2\)（\(0<t\le\pi\)），
\(\|D_N\|_1\) 与
\[
 I_N=\int_0^{\lambda\pi}\frac{|\sin u|}{u}\,du
\]
上下相差固定正倍数。在初始块 \([0,\pi]\)，积分为固定正数且有限。
对完整块 \([j\pi,(j+1)\pi]\)、\(j\ge1\)，有
\[
 \frac{2}{(j+1)\pi}\le
 \int_{j\pi}^{(j+1)\pi}\frac{|\sin u|}{u}\,du
 \le\frac{2}{j\pi}.
\]
当 \(N\ge2\) 时，完整块包括 \(j=1,\ldots,N-1\)，最后半块的积分
不超过 \(1/N\)。调和和的估计给出
\[
 c\log(N+2)\le\|D_N\|_1\le C\log(N+2)\qquad(N\ge0),
\]
其中 \(c,C>0\) 与 \(N\) 无关；\(N=0,1\) 由调整常数包含进去。

固定一点 \(x\)，在 Banach 空间 \(C(\mathbb T)\) 上令
\(T_Nf=S_Nf(x)\)。卷积公式给
\(\|T_N\|\le(2\pi)^{-1}\|D_N\|_1\)。反向取连续周期函数
\[
 f_\eta(s)=\frac{D_N(x-s)}{\sqrt{D_N(x-s)^2+\eta^2}},\qquad \eta>0.
\]
其一致范数不超过一，且由控制收敛，
\[
 T_Nf_\eta=\frac1{2\pi}\int_{-\pi}^{\pi}
 \frac{D_N(t)^2}{\sqrt{D_N(t)^2+\eta^2}}\,dt
 \longrightarrow\frac1{2\pi}\|D_N\|_1.
\]
所以两者相等，定点泛函的范数无共同上界。一致有界原理的逆否命题
便给出某个 \(f\in C(\mathbb T)\)，使 \(\sup_N|S_Nf(x)|=\infty\)。
特别地，连续性并不足以保证 Fourier 部分和在每一点收敛。

\section{局部收敛判据}

点 \(x\) 附近的左右极限若存在，记
\[
 s_x=\frac{f(x+)+f(x-)}2.
\]
由 Dirichlet 核的偶性和积分为 \(2\pi\)，公式
\eqref{eq:fourier-dirichlet-convolution}可写成
\begin{equation}\label{eq:fourier-symmetric-difference}
 S_Nf(x)-s_x
 =\frac1{2\pi}\int_0^\pi
 \frac{f(x+t)+f(x-t)-2s_x}{2\sin(t/2)}
 2\sin\bigl((N+\tfrac12)t\bigr)\,dt.
\end{equation}

\begin{lemma}[圆周上的 Riemann--Lebesgue 引理]\label{lem:circle-riemann-lebesgue}
若 \(g\in L^1(\mathbb T)\)，则
\(\int_{-\pi}^{\pi}g(t)e^{-ikt}\,dt\to0\)，其中整数频率满足 \(|k|\to\infty\)。
\end{lemma}
\begin{proof}
平移在 \(L^1(\mathbb T)\) 中连续。对 \(k\ne0\) 取
\(h=\pi/k\)，周期换元给
\[
 2\int g(t)e^{-ikt}\,dt
 =\int\bigl(g(t)-g(t-h)\bigr)e^{-ikt}\,dt.
\]
故左侧绝对值不超过 \(\|g-\tau_hg\|_1\)，而 \(h\to0\)。平移连续性可先
对连续周期函数由一致连续性证明，再用连续函数在 \(L^1(\mathbb T)\) 中的
稠密性推广。
\end{proof}

\begin{theorem}[Dini 判据]\label{thm:fourier-dini-test}
若 \(f\in L^1(\mathbb T)\)，点 \(x\) 的左右极限存在，并且对某个
\(\delta>0\)，
\[
 t\longmapsto
 \frac{f(x+t)+f(x-t)-f(x+)-f(x-)}{t}
\]
在 \((0,\delta)\) 可积，则 \(S_Nf(x)\to s_x\)。
\end{theorem}
\begin{proof}
先缩小 \(\delta\)，使 \(0<\delta<\pi\)。
在 \((0,\delta)\) 上，\(t/(2\sin(t/2))\) 有界，故
\eqref{eq:fourier-symmetric-difference}的被积振幅可积。区间
\([\delta,\pi]\) 上分母离零有界，而 \(f(x+t)+f(x-t)\) 可积。于是整个振幅
属于 \(L^1(0,\pi)\)，引理应用于频率 \(N+\tfrac12\) 即得结论；半整数频率
也可把正弦拆成整数频率乘以固定的 \(\sin(t/2),\cos(t/2)\)。
\end{proof}

本章说周期函数分段 \(C^1\)，是指一个周期可作有限分割，使每个开片上的
限制延拓为相应闭片上的 \(C^1\) 函数；分割点处原函数的取值可以另定。
因此单侧邻域的导数有界，中值定理给
\(|f(x\pm t)-f(x\pm)|\le Ct\)，Dini 条件成立。
仅在开片内部 \(C^1\) 并不保证端点附近的导数有界。
前面构造的连续函数发散例子则说明，一致逼近的三角多项式不能任意换成
原函数自己的 Fourier 部分和。

\begin{theorem}[Dirichlet--Jordan 定理]\label{thm:dirichlet-jordan}
若 \(f\) 在圆周上有界变差，则对每个 \(x\)
\[
 S_Nf(x)\longrightarrow\frac{f(x+)+f(x-)}2.
\]
特别地，在连续点收敛到 \(f(x)\)，在跳跃点收敛到左右极限的平均值。
\end{theorem}
\begin{proof}
先建立与频率无关的核积分界。函数
\[
 r(t)=\frac1{\sin(t/2)}-\frac2t
\]
在 \(t=0\) 可连续补值为零，因而在 \([0,\pi]\) 有界。
同时 \(\int_a^b\sin u/u\,du\) 对所有 \(0\le a\le b\) 一致有界：
在 \([0,1]\) 用 \(|\sin u/u|\le1\)，在 \([1,b]\) 分部积分后，
边界项和 \(\int_1^b u^{-2}du\) 都有固定界；一般子区间是两个初始
积分之差。因此存在绝对常数 \(C\)，使
\[
 \left|\int_a^b\frac{\sin(\lambda t)}{\sin(t/2)}\,dt\right|
 \le C\qquad(0\le a\le b\le\pi,\ \lambda>0).
\]
这里 \(2/t\) 项换元为正弦积分，余项由
\(\int_0^\pi|r(t)|\,dt\) 控制。

设 \(u\) 在 \([0,\pi]\) 上右连续、有界变差且 \(u(0)=0\)。
由\chapterref{chap:signed-measures}的 Stieltjes 构造，有有限符号测度
\(\nu\) 满足 \(u(t)=\nu((0,t])\)；复值时分别处理实部和虚部。
对固定 \(\lambda>0\)，核在零处有有限极限 \(2\lambda\)，所以有限区间上
有界。于是下式的换序绝对可积，Fubini 给出
\[
 \int_0^\delta u(t)\frac{\sin(\lambda t)}{\sin(t/2)}\,dt
 =\int_{(0,\delta]}
       \left(\int_s^\delta\frac{\sin(\lambda t)}{\sin(t/2)}\,dt\right)d\nu(s).
\]
其绝对值不超过 \(C|\nu|((0,\delta])\)。因
\((0,\delta]\downarrow\varnothing\)，有限变差测度的连续性使此界随
\(\delta\downarrow0\) 趋零，而且与 \(\lambda\) 无关。
固定 \(\delta>0\) 后，\(u(t)/\sin(t/2)\) 在 \([\delta,\pi]\) 可积；
取 \(\lambda=N+\tfrac12\)，由 Riemann--Lebesgue 引理（拆开半整数频率）
知远区积分趋零。先选小 \(\delta\)，再令 \(N\to\infty\)，即得
\[
 \int_0^\pi u(t)\frac{\sin((N+\tfrac12)t)}{\sin(t/2)}\,dt\longrightarrow0.
\]

对固定 \(x\)，右差 \(f(x+t)-f(x+)\) 和左差 \(f(x-t)-f(x-)\)
在 \(t\downarrow0\) 时均趋零，且都是 BV 函数。它们未必关于 \(t\)
右连续，但可各取右连续代表，在 \(t=0\) 取零：BV 函数只有至多可数个
间断点，这种修改不改变上述积分，也不改变零处的右极限或有界变差性质。
分别应用刚证结论，再代入\eqref{eq:fourier-symmetric-difference}，即得定理。
\end{proof}

\section{跳跃与 Gibbs 现象}

Dirichlet--Jordan 定理只描述固定点的极限。若观察点随 \(N\) 向跳跃点靠近，
部分和会在越来越窄的区间内出现不消失的超调。以周期阶跃函数
\[
 f(x)=\begin{cases}-1,&-\pi<x<0,\\1,&0<x<\pi\end{cases}
\]
为例，直接计算得
\[
 S_Nf(x)=\frac4\pi\sum_{\substack{1\le k\le N\\k\ {\rm odd}}}
 \frac{\sin(kx)}k.
\]
令 \(x=y/N\)，把和视为 Riemann 和，可得局部极限
\[
 S_Nf(y/N)\longrightarrow \frac2\pi\int_0^y\frac{\sin u}{u}\,du.
\]
记右侧的局部极限函数为 \(G(y)\)。由
\(G'(y)=2\sin y/(\pi y)\)，它在正半轴的第一个极大点恰为 \(y=\pi\)。
在观察点 \(x=\pi/N\)，超过右极限 \(1\) 的绝对量趋于
\[
 A=\frac2\pi\int_0^\pi\frac{\sin u}{u}\,du-1\approx0.17898.
\]
跳跃高度为二，所以相对于跳跃高度的超调比例是 \(A/2\approx0.08949\)，
即约 \(8.949\%\)。这区分了绝对超调量和通常所说的 Gibbs 比例。
观察点趋向跳跃点，但超调高度不趋零；这不与 \(L^2\) 收敛矛盾，
因为积分收敛不控制逐点误差。该函数的 \(L^2\) 收敛将在下一章由正交展开证明，
不能只凭超调区域看起来很窄就断言整体平方误差趋零。

\section{Fejér 平均}

定义部分和的 Cesàro 平均
\[
 \sigma_Nf=\frac1{N+1}\sum_{n=0}^NS_nf.
\]
交换有限求和得
\[
 \sigma_Nf(x)=\frac1{2\pi}\int_{-\pi}^{\pi}f(x-t)K_N(t)\,dt,
\]
其中 Fejér 核
\begin{equation}\label{eq:fejer-kernel}
 K_N(t)=\frac1{N+1}\sum_{n=0}^ND_n(t)
 =\frac1{N+1}\left(\frac{\sin((N+1)t/2)}{\sin(t/2)}\right)^2.
\end{equation}
在 \(t\in2\pi\mathbb Z\) 处按可去值定义 \(K_N(t)=N+1\)。它满足
\[
 K_N\ge0,\qquad \frac1{2\pi}\int_{-\pi}^{\pi}K_N=1,\qquad
 \sup_{\delta\le|t|\le\pi}K_N(t)\longrightarrow0
\]
对每个 \(0<\delta\le\pi\) 成立。三条性质使 \(K_N\) 成为圆周上的逼近恒等。

\begin{theorem}[Fejér 定理]\label{thm:fejer-uniform}
若 \(f\in C(\mathbb T)\)，则 \(\sigma_Nf\to f\) 一致收敛。更一般地，若
\(f\in L^1(\mathbb T)\)，则在每个 Lebesgue 点 \(x\) 有
\(\sigma_Nf(x)\to f(x)\)。
\end{theorem}
\begin{proof}
连续情形中
\[
 |\sigma_Nf(x)-f(x)|
 \le\frac1{2\pi}\int_{-\pi}^{\pi}|f(x-t)-f(x)|K_N(t)\,dt.
\]
给定 \(\varepsilon>0\)，一致连续性给 \(\delta>0\)，使
\(|t|<\delta\) 时差小于 \(\varepsilon\)，这部分积分至多
\(\varepsilon\)。外部积分由 \(2\|f\|_\infty\) 乘以核在
\(\delta\le|t|\le\pi\) 上的质量控制，后一质量趋零；估计与 \(x\) 无关。

若 \(x\) 为 Lebesgue 点，令 \(F(t)=|f(x-t)-f(x)|\)，并记
\(H(s)=\int_{|t|\le s}F(t)\,dt\)。于是 \(H(s)/s\to0\)。
核公式与 \(|\sin(t/2)|\ge |t|/\pi\) 给
\[
 K_N(t)\le C\min\{N+1,((N+1)t^2)^{-1}\}.
\]
其中 \(C\) 为绝对常数；第一项也可由
\(|\sum_{k=0}^Ne^{ikt}|\le N+1\) 得到。
给定 \(\varepsilon>0\)，取 \(0<\delta<\pi\) 使
\(H(s)\le\varepsilon s\)（\(0<s\le\delta\)），再取 \(N\) 使
\(\rho=(N+1)^{-1}<\delta\)。中心区间的积分满足
\[
 \int_{|t|\le\rho}F(t)K_N(t)\,dt
 \le C(N+1)H(\rho)\le C\varepsilon.
\]
在 \(\rho<|t|\le\delta\) 上，按
\(a_j=2^j\rho<\delta\)、\(b_j=\min(2a_j,\delta)\) 分成环带。
每一带（包括末端截断的一带）均满足
\[
 \int_{a_j<|t|\le b_j}F(t)K_N(t)\,dt
 \le\frac{C H(b_j)}{(N+1)a_j^2}
 \le 2C\varepsilon\,2^{-j}.
\]
求和后近区总量不超过 \(5C\varepsilon\)。远区则有
\[
 \int_{\delta<|t|\le\pi}F(t)K_N(t)\,dt
 \le\frac{C}{(N+1)\delta^2}\int_{-\pi}^{\pi}F(t)\,dt\longrightarrow0.
\]
因为 \(f\in L^1\) 且 Lebesgue 点的 \(f(x)\) 有限，最后的积分有限。
乘回 \(1/(2\pi)\)，先令 \(N\to\infty\)，再令
\(\varepsilon\downarrow0\)，便得结论。
\end{proof}

连续情形立即证明三角多项式在 \(C(\mathbb T)\) 中一致稠密，因为
\(\sigma_Nf\) 本身是次数不超过 \(N\) 的三角多项式。这里使用的是经过平均的
部分和；不能把结论中的 \(\sigma_Nf\) 替换为 \(S_Nf\)。

\section*{习题}
\addcontentsline{toc}{section}{习题}
\subsection*{A组}
\begin{enumerate}[label=\textbf{\thechapter.\arabic*.}]
\exerciseitem{fourier-circle-real-coefficients}把复 Fourier 系数与正弦、余弦系数互相换算，并证明实值函数的共轭对称性。
\exerciseitem{fourier-dirichlet-kernel-formula}由有限等比和推导 Dirichlet 核公式及其总质量。
\exerciseitem{fourier-translation-coefficients}若 \(g(x)=f(x-a)\)，证明 \(\widehat g(k)=e^{-ika}\widehat f(k)\)。
\exerciseitem{fourier-step-series}计算正文周期阶跃函数的 Fourier 系数并写出部分和。
\end{enumerate}
\subsection*{B组}
\begin{enumerate}[label=\textbf{\thechapter.\arabic*.},resume]
\exerciseitem{fourier-dini-piecewise-c-one}证明分段 \(C^1\) 周期函数满足 Dini 判据，并写出跳跃点的极限。
\exerciseitem{fourier-localization}证明若两个 \(L^1\) 周期函数在 \(x\) 的某邻域相等，则它们 Fourier 部分和在 \(x\) 处之差趋零。
\exerciseitem{fourier-fejer-three-properties}逐项验证 Fejér 核的非负性、总质量和远离原点的一致趋零。
\exerciseitem{fourier-fejer-positive-contraction}证明 \(\|\sigma_Nf\|_p\le\|f\|_p\)，其中 \(1\le p\le\infty\)。
\exerciseitem{fourier-fejer-lp}证明对 \(1\le p<\infty\)，\(\sigma_Nf\to f\) 于 \(L^p(\mathbb T)\)。
\end{enumerate}
\subsection*{C组}
\begin{enumerate}[label=\textbf{\thechapter.\arabic*.},resume]
\exerciseitem{fourier-uniqueness-lone}若 \(f\in L^1(\mathbb T)\) 的全部 Fourier 系数为零，利用 Fejér 平均证明 \(f=0\) 几乎处处。
\optionalexerciseitem{fourier-dirichlet-lone-growth}证明 \(\|D_N\|_{L^1(-\pi,\pi)}\) 的量级为 \(\log(N+2)\)。
\exerciseitem{fourier-gibbs-scaling}补全阶跃函数在 \(x=y/N\) 下的 Riemann 和极限，并解释为何超调不妨碍 \(L^2\) 收敛。
\end{enumerate}
